\begin{textsample}{NEG Dim 3 – 2023\_c5 – Score 3.00 – 2023\_c5/enade\_2023\_fisio\_8b895e4d42.txt}  \label{ex:f3_neg_225}
CAPÍTULO 4 DISTRIBUIÇÃO DE CONCEITOS A finalidade deste capítulo é apresentar um panorama nacional sobre a distribuição do Conceito Enade 2023.
Para traçar esse panorama, recorre -se a tabelas e gráficos, que são acompanhados de comentários e análises que contribuem para a contextualização dos dados neles expostos.
Para visualizar o panorama do Enade 2023, em nível nacional, na primeira seção deste capítulo, os dados sobre o número de cursos participantes são organizados por grande região.
Na segunda, o mapeamento do conceito do exame é feito por categoria administrativa das IES e modalidade de oferta dos cursos, de maneira cruzada às grandes regiões.
Na última seção, a lógica de apresentação dos dados é feita por organização acadêmica e, também, por grande região de funcionamento dos cursos, no país como um todo. 4.1 PANORAMA NACIONAL DA DISTRIBUIÇÃO DOS CONCEITOS A lógica adotada na apresentação dos dados, nesta seção e nas duas seguintes, pressupõe o registro de duas informações de natureza reiterativa.
O primeiro diz respeito aos valores possíveis do Enade, que variam em níveis de 1 a 5, sendo que, à medida que esses valores aumentam, o curso é melhor avaliado, e ao fato de que aqueles cursos que tiveram somente um ou nenhum concluinte participante não obtêm o Conceito Enade, ficando sem conceito ( SC ), como ressaltado na apresentação deste relatório.
O segundo registro é que, caso haja diferença entre os cursos tabulados neste capítulo e no capítulo 1, ela corresponderá aos cursos que não foram avaliados, ou seja, àqueles, em princípio, sem estudantes concluintes que tivessem sido inscritos para a prova.
Na Tabela 4.1, são apresentadas a quantidade e a distribuição de cursos de Fisioterapia participantes do Enade 2023, por faixa de conceito e grande região.
A partir dos dados da Tabela 4.1, nota -se que, dos 711 cursos participantes, 319 ( 44,9 \% ) classificaram -se com conceito 3, o valor modal.
Esse foi também o conceito modal nas regiões : Nordeste ( 41,9 \% ), Sudeste ( 49,1 \% ), Sul ( 45,7 \% ).
Nas regiões Centro-Oeste e Norte, o conceito modal foi o 2 ( 40,0 \% e 45,5 \%, respectivamente ).
Em nível nacional, o conceito 2 foi o segundo mais frequente ( 27,0 \%, correspondendo a 192 cursos ), e o conceito 4 foi o terceiro ( 19,1 \%, correspondendo a 136 cursos ).
Houve, ainda, 38 ( 5,3 \% ) cursos que receberam o conceito 5 e 14 ( 2,0 \% ) cursos que receberam o conceito 1.
Dos 711 cursos de Fisioterapia, 12 ( 1,7 \% ) ficaram sem conceito ( SC ).
A região Norte participou com 44 cursos, o que corresponde a 6,2 \% do total nacional.
Desses, 20 cursos receberam o conceito 2, o que equivale a 45,5 \% do total regional, o conceito modal na região, conforme mencionado.
A 19 cursos ( 43,2 \% ) atribuiu -se o conceito 3 ; a três cursos ( 6,8 \% ), o conceito 4 ; e a um curso ( 2,3 \% ), o conceito 5.
Ainda nessa região, um curso ( 2,3 \% ) obteve conceito 1 e nenhum curso ficou sem conceito ( SC ).
A região Nordeste participou com 179 cursos, correspondentes a 25,2 \% do total nacional.
Nessa região, 75 cursos ( 41,9 \% em termos regionais ) obtiveram conceito 3, o conceito modal para a região, como já comentado.
Os conceitos 2 e 4 foram atribuídos, respectivamente, a 65 e a 17 cursos ( 36,3 \% e 9,5 \%, respectivamente ).
Ainda nessa região, 12 cursos ( 6,7 \% ) foram avaliados com o conceito 5 e três cursos ( 1,7 \% ) foram avaliados com o conceito 1.
Outros sete cursos ( 3,9 \% ) ficaram sem conceito ( SC ).
Para os 291 cursos participantes da região Sudeste ( 40,9 \% do total nacional ), o conceito 3 foi o conceito modal, como já anteriormente assinalado, tendo sido obtido por 143 cursos ( 49,1 \% ).
O conceito 4 foi atribuído a 65 cursos ( 22,3 \% ) e o conceito 2, a 60 cursos ( 20,6 \% ).
Receberam os conceitos 5 e 1, respectivamente, 14 cursos ( 4,8 \% ) e cinco cursos ( 1,7 \% ).
Nessa região, quatro cursos ( 1,4 \% ) ficaram sem conceito ( SC ).
Os 127 cursos da região Sul corresponderam a 17,9 \% do total nacional.
Como já indicado, houve predominância do conceito 3 ( conceito modal ), atribuído a 45,7 \% dos cursos da região, ou seja, a 58 dos 127 participantes.
O conceito 4 foi atribuído a 39 cursos ( 30,7 \% ) ; o conceito 2, 19 cursos ( 15,0 \% ) ; e o conceito 5, a dez cursos ( 7,9 \% ).
Na região Sul, nenhum curso obteve conceito 1 e um curso ( 0,8 \% ) ficou sem conceito ( SC ).
Já dos 70 cursos participantes na região Centro-Oeste ( 9,8 \% do total nacional ), 28 ( 40,0 \% em termos regionais ) receberam conceito 2, o conceito modal, como já destacado.
Outros 24 cursos ( 34,3 \% ) obtiveram o conceito 3, 12 cursos ( 17,1 \% ) obtiveram o conceito 4 e cinco cursos ( 7,1 \% ), o conceito 1.
Nessa região, um curso ( 1,4 \% ) obteve o conceito 5 e nenhum curso ficou sem conceito ( SC ).
Em um gráfico de distribuição cumulativa, em cada valor de abscissa ( eixo x ), o valor apresentado no eixo das ordenadas ( eixo y ) é o valor acumulado da variável ( Conceito Enade ).
Em linhas gerais, em um gráfico desta mesma tipologia, poligonais mais à esquerda representam, no caso dos dados apresentados nesta seção, as grandes regiões com a pior distribuição, e poligonais mais à direita, as grandes regiões com os melhores conceitos.
Considerando essas diretrizes, no Gráfico 4.1, é apresentada a distribuição cumulativa do Conceito Enade, segundo a grande região.
No Gráfico 4.1, observando -se a linha vertical pontilhada, que ajuda a visualizar o percentual de conceitos acumulados até o valor 3, ou seja, conceitos 1, 2 e 3, pode -se notar, por exemplo, que a região Norte ( linha azul ) apresenta 90,9 \% dos cursos com conceito 3 ou abaixo desse valor.
A região Sul ( linha roxa ), por outro lado, apresenta 61,1 \% dos cursos na mesma situação : com conceito 3 ou abaixo.
Isso significa que a região Norte apresenta os cursos com conceitos mais baixos, ao passo que a região Sul apresenta a distribuição com valores maiores.
Os cursos em IES situadas na região Sudeste ( linha vermelha ) apresentam uma situação intermediária. 4.2 CONCEITOS POR CATEGORIA ADMINISTRATIVA, MODALIDADE DE OFERTA DOS CURSOS E GRANDE REGIÃO Na Tabela 4.2, é apresentada a distribuição dos cursos de Fisioterapia participantes do Enade 2023, por categoria administrativa e modalidade de oferta, de acordo com os conceitos por eles alcançados, segundo as grandes regiões brasileiras.
De acordo com as informações da Tabela 4.2, dos 711 cursos participantes, 61 ( 8,6 \% ) eram ministrados em instituições públicas e 650 ( 91,4 \% ), em instituições privadas.
Quanto à modalidade de oferta, 683 ( 96,1 \% ) cursos eram presenciais e 28 ( 3,9 \% ), a distância.
Conforme informações da mesma tabela, em termos nacionais, entre instituições públicas e privadas, observa -se que, dos 38 cursos avaliados com conceito 5, 25 foram oferecidos em IES públicas, e 13, em IES privadas.
Dos 650 cursos participantes de IES privadas, o conceito 3 foi o valor modal, atribuído a 314 ( 48,3 \% ) cursos.
Nessa mesma categoria, 187 cursos ( 28,8 \% ) foram avaliados com o conceito 2, 110 cursos ( 16,9 \% ) foram avaliados com o conceito 4, 14 ( 2,2 \% ) foram avaliados com o conceito 1 e 13 ( 2,0 \% ), com o conceito 5.
Ademais, 12 cursos ( 1,8 \% ) ficaram sem conceito ( SC ).
Na rede pública, o conceito modal foi o 4, atribuído a 26 cursos, correspondentes a 42,6 \% dos 61 da categoria.
Dos demais cursos participantes, 25 ( 41,0 \% ) obtiveram o conceito 5, cinco ( 8,2 \% ) receberam o conceito 3 e outros cinco ( 8,2 \% ), o conceito 2.
Nenhum curso recebeu o conceito 1 ou ficou sem conceito ( SC ) nessa categoria administrativa.
Quanto à modalidade de oferta, dos 683 cursos presenciais, 313 ( 45,8 \% ) receberam o conceito 3, valor modal nessa categoria.
Dos demais cursos, 175 ( 25,6 \% ) receberam o conceito 2 ; 135 ( 19,8 \% ) receberam o conceito 4 ; 37 ( 5,4 \% ), o conceito 5 ; e 14 ( 2,0 \% ), o conceito 1.
Nessa modalidade de oferta, nove cursos ( 1,3 \% ) ficaram sem conceito ( SC ).
Para os 28 cursos participantes na modalidade a distância, o conceito 2 foi o valor modal, atribuído a 17 ( 60,7 \% ) deles.
Outros seis cursos ( 21,4 \% ) receberam o conceito 3.
Os conceitos 4 e 5 foram atribuídos a um curso cada ( 3,6 \% ).
Nenhum curso recebeu o conceito 1 e três ( 10,7 \% ) ficaram sem conceito ( SC ).
Na análise por região, observa -se, que, na região Norte, as instituições públicas participaram com cinco cursos ( 11,4 \% total regional ), dos quais nenhum ficou sem conceito ( SC ) ou recebeu o conceito 1.
Os conceitos 2 e 4 apresentaram os maiores quantitativos, com dois cursos em cada, e um curso obteve o conceito 5 ( 20,0 \% ).
Já as instituições privadas participaram com 39 cursos ( 88,6 \% do total regional ), dos quais nenhum ficou sem conceito ( SC ) ou alcançou o conceito 5.
O conceito modal nessa categoria na região Norte foi o 3, atribuído a 19 cursos, o que corresponde a 48,7 \%.
O conceito 2 foi atribuído a 18 cursos, e os conceitos 1 e 4 foram atribuídos a um curso cada.
Dos cursos oferecidos na região Norte, 43 eram presenciais, sendo que nenhum ficou sem conceito ( SC ).
Nessa modalidade, 19 obtiveram o conceito 3 ; outros 19, o conceito 2 ; três, o conceito 4 ; um, o conceito 5 ; e um, o conceito 1.
Um curso foi registrado na modalidade a distância, que recebeu o conceito 2.
Na região Nordeste, a rede privada concentra 163 dos 179 cursos participantes, o equivalente a 91,1 \% do total da região.
Entre esses, 75 cursos foram avaliados com conceito 3, o conceito modal ; 64 cursos obtiveram o conceito 2 ; 12 cursos, o conceito 4 ; três cursos, o conceito 1 ; e dois cursos, o conceito 5.
Nessa categoria administrativa e grande região, sete cursos ficaram sem conceito ( SC ).
As instituições públicas da região Nordeste participaram com 16 cursos ( 8,9 \% do total da região ).
Desses, dez ( 62,5 \% ) alcançaram conceito 5, o valor modal, seguidos por cinco cursos que receberam o conceito 4.
O conceito 2 foi atribuído a um curso.
Nenhum curso oobteve os conceitos 1 e 3 ou ficou sem conceito ( SC ).
Dos cursos oferecidos na região Nordeste, 174 eram presenciais : seis ficaram sem conceito ( SC ) ; três obtiveram o conceito 1 ; 61, o conceito 2 ; 75, o conceito 3 ( valor modal ) ; 17, o conceito 4 ; e 12, o conceito 5.
Cinco cursos eram na modalidade a distância, sendo que quatro receberam o conceito modal 2 e um ficou sem conceito ( SC ).
Na região Sudeste, a proporção de cursos da rede privada, 92,1 \%, correspondeu a 268 dos 291 cursos participantes.
Nessa combinação de categoria administrativa e grande região, o conceito modal foi o 3, atribuído a 140 cursos ; e quatro cursos ficaram sem conceito ( SC ).
Os demais foram avaliados com os seguintes conceitos : 2 ( 59 cursos ), 4 ( 54 cursos ), 5 ( seis cursos ) e 1 ( cinco cursos ).
Dos 23 cursos oferecidos por instituições públicas na região Sudeste, oito obtiveram o conceito 5 ; 11, o conceito 4 ; três, o conceito 3 ; e um obteve o conceito 2.
Nenhum curso obteve o conceito 1 ou ficou sem conceito ( SC ).
Ainda na região Sudeste, grande parcela dos cursos ( 96,9 \% ) foi oferecida na modalidade presencial : 282 dos 291 cursos.
Para esses, o conceito modal foi o 3, com 140 cursos.
Os demais cursos presenciais da região foram alocados nos conceitos 4 ( 65 cursos ), 2 ( 56 cursos ), 5 ( 14 cursos ) e 1 ( cinco cursos ).
Dois cursos presenciais ficaram sem conceito ( SC ).
Na modalidade a distância, nove cursos foram oferecidos, sendo que dois ficaram sem conceito ( SC ), quatro alcançaram o conceito 2 e três cursos, o conceito 3.
As instituições privadas concentraram 116 dos 127 cursos participantes da região Sul, o que representa 91,3 \% do total regional.
Desses, 57 obtiveram conceito 3, o conceito modal.
Nessa combinação de categoria administrativa e grande região, 35 receberam o conceito 4 ; 19 receberam o conceito 2 ; quatro, o conceito 5 ; e nenhum obteve o conceito 1.
Um curso ficou sem conceito ( SC ).
As instituições públicas na região Sul participaram com 11 cursos ( 8,7 \% dos da região ), aos quais foram atribuídos os conceitos 3 ( um curso ), 4 ( quatro cursos ) e 5 ( seis cursos ).
Nenhum curso da região Sul oferecido por IES públicas ficou sem conceito ( SC ).
Na região Sul, 115 cursos foram oferecidos na modalidade presencial.
Para esses, o valor modal foi o conceito 3, com 55 cursos.
Os demais cursos presenciais da região foram alocados nos conceitos 4 ( 38 cursos ), 2 ( 12 cursos ) e 5 ( nove cursos ).
Um curso ficou sem conceito ( SC ).
Na modalidade a distância, foram oferecidos 12 cursos, sendo sete cursos com conceito 2 ( valor modal ), três cursos com conceito 3, um curso com conceito 4 e um curso com conceito 5.
Nenhum curso ficou sem conceito ( SC ) nessa região.
Na região Centro-Oeste, 64 dos 70 cursos participantes eram de IES privadas ( 91,4 \% em termos regionais ).
Desses, 27 cursos receberam conceito 2, o valor modal.
Dos demais cursos, 23 receberam o conceito 3, oito receberam o conceito 4, cinco receberam o conceito 1 e um recebeu o conceito 5.
Nenhum ficou sem conceito ( SC ).
Quanto aos seis cursos oferecidos por instituições públicas nessa região ( 8,6 \% do total regional ), um ficou com conceito 2 ; um, com conceito 3 ; e quatro, com conceito 4.
Da mesma forma, nenhum ficou sem conceito ( SC ) nessa categoria.
A maioria dos cursos ofertados na região Centro-Oeste eram na modalidade presencial : 69 ( 98,6 \% ).
Para esses, o valor modal foi o conceito 2, com 27 cursos ( 39,1 \% ).
Os demais foram alocados nos conceitos 3 ( 24 cursos, 34,8 \% ), 4 ( 12 cursos, 17,4 \% ), 1 ( cinco cursos, 7,2 \% ) e 5 ( um curso, 1,4 \% ).
Nenhum curso presencial da região ficou sem conceito ( SC ).
Na modalidade a distância, foi oferecido somente um curso, que alcançou o conceito 2.
A título de síntese, no Gráfico 4.2, apresenta -se a distribuição do Conceito Enade, segundo a categoria administrativa da IES.
Nesse gráfico, como dito anteriormente, poligonais mais à esquerda representam Conceito Enade mais baixos, e poligonais mais à direita, Conceito Enade mais altos.
É importante observar que, para cada poligonal presente nesse gráfico, a porcentagem constante no eixo y refere -se ao universo de cursos oferecidos por instituições públicas ( poligonal azul ) ou privadas ( poligonal verde ), respectivamente.
Os cursos em IES públicas apresentam uma poligonal mais à direita do que os em IES privadas e, consequentemente, uma distribuição de Conceito Enade com valores maiores. 4.3 CONCEITOS POR ORGANIZAÇÃO ACADÊMICA E GRANDE REGIÃO Na Tabela 4.3, encontra -se a distribuição dos conceitos atribuídos aos cursos participantes do Enade 2023, na área de Fisioterapia, por organização acadêmica, segundo as grandes regiões brasileiras.
Como mostram os dados da Tabela 4.3, dos 711 cursos de Fisioterapia participantes, 214 eram oferecidos em universidades, 287 em centros universitários, 209 em faculdades e um curso foi oferecido por CEFET/IF.
Essa distribuição corresponde a, respectivamente, 30,1 \%, 40,4 \%, 29,4 \% e 0,1 \% dos cursos.
De acordo com os dados apresentados na mesma tabela, de todos os 136 cursos avaliados com o conceito 4, 72 eram vinculados a universidades.
Os cursos vinculados a esse tipo de organização acadêmica tiveram o conceito 3 como conceito modal, atribuído a 75 cursos ( 35,0 \% ).
Os demais cursos ofertados por universidades receberam os conceitos na seguinte proporção : 34 cursos receberam conceito 2 ; 29 cursos, conceito 5 ; e um curso, conceito 1.
Três cursos ficaram sem conceito ( SC ).
Dos cursos vinculados a centros universitários, o conceito modal também foi o 3, atribuído a 154 cursos ( 53,7 \% ).
Cinco foi o quantitativo de cursos vinculados a esse tipo de organização acadêmica que ficaram sem conceito ( SC ) ou alcançaram os conceitos 1 ou 5.
Dos demais, 71 receberam o conceito 2 e 47, o conceito 4.
Entre os 209 cursos mantidos por faculdades, 90 ( 43,1 \% ) receberam conceito 3, o conceito modal.
O segundo conceito mais frequentemente atribuído aos cursos oferecidos por essa categoria de IES foi o conceito 2, recebido por 87 cursos.
Na sequência, 16 cursos obtiveram conceito 4 ; oito cursos, o conceito 1 ; e quatro cursos, o conceito 5.
Dos cursos oferecidos em faculdades, quatro ficaram sem conceito ( SC ).
Somente um curso foi oferecido por CEFET/IF, tendo alcançado o conceito 4.
Considerando -se, separadamente, as regiões brasileiras, verifica -se que, na região Norte, as universidades participaram com nove cursos avaliados, aos quais foram atribuídos os seguintes conceitos : 2, o conceito modal, a quatro cursos ; 3, a dois cursos ; 4, a dois cursos ; e 5, a um curso.
Nenhum curso ficou sem conceito ( SC ).
Os centros universitários da região Norte foram representados por 18 cursos, dos quais nenhum obteve os conceitos 1 ou 5 ou ficou sem conceito ( SC ).
Dos demais, oito cursos alcançaram o conceito 2 ; nove cursos, conceito 3, o valor modal ; e um curso, o conceito 4.
As faculdades participaram com 17 cursos na região Norte, os quais receberam os conceitos 1 ( um curso ), 2 ( oito cursos ) e 3 ( oito cursos ).
Nenhum curso ficou sem conceito ( SC ) ou alcançou os conceitos 4 ou 5.
Por fim, não houve registros de cursos ofertados por CEFET/IF na região Norte.
Na região Nordeste, as universidades participaram com 29 dos 179 cursos da área de Fisioterapia oferecidos.
O conceito modal foi 5, atribuído a dez cursos.
Os conceitos 2, 3 e 4 foram atribuídos a seis cursos cada.
Um curso ficou sem conceito ( SC ) e nenhum curso obteve o conceito 1 Os centros universitários contaram com 80 cursos participantes na região Nordeste, 41 dos quais receberam o conceito modal 3.
Para os demais cursos, foram atribuídos os conceitos 2 ( 27 cursos ), 4 ( sete cursos ), 5 ( dois cursos ) e 1 ( um curso ).
Dois cursos oferecidos por centros universitários no Nordeste ficaram sem conceito ( SC ).
As faculdades foram representadas por 70 cursos na região Nordeste, os quais receberam os conceitos 1 ( dois cursos ), 2 ( 32 cursos ), 3 ( 28 cursos ) e 4 ( quatro cursos ).
Quatro cursos ficaram sem conceito ( SC ).
Da mesma forma que a região Norte, não houve registros de cursos ofertados por CEFET/IF na região Nordeste.
Na região Sudeste, as universidades concentraram 107 dos 291 cursos de Fisioterapia oferecidos.
O conceito modal nesse tipo de organização acadêmica foi o 3, atribuído a 44 cursos.
Os demais receberam os conceitos 4 ( 30 cursos ), 2 ( 19 cursos ), 5 ( 11 cursos ) e 1 ( um curso ).
Dois cursos ficaram sem conceito ( SC ).
Os centros universitários participaram com 115 cursos na região Sudeste, dos quais 63 obtiveram o conceito modal 3.
O conceito 4 foi atribuído a 24 cursos ; o conceito 2, a 22 cursos ; o conceito 1, a três cursos ; e o conceito 5, a um curso.
Dois cursos ficaram sem conceito ( SC ) nessa combinação de organização acadêmica e região.
Já as faculdades foram representadas por 68 cursos na região Sudeste, dos quais nenhum ficou sem conceito ( SC ).
O conceito modal 3 foi atribuído a 36 cursos.
Aos demais cursos foram atribuídos os conceitos 2 ( 19 cursos ), 4 ( dez cursos ), 5 ( dois cursos ) e 1 ( um curso ).
Os CEFET/IF participaram com um curso na região Sudeste, ao qual foi atribuído o conceito 4.
Dos 127 cursos da região Sul, 55 eram mantidos por universidades, para os quais o conceito modal foi 4, atribuído a 28 cursos.
Os demais cursos receberam os conceitos 3 ( 17 cursos ), 5 ( sete cursos ) e 2 ( três cursos ).
Em IES com esse tipo de organização acadêmica, nenhum curso da região Sul ficou sem conceito ( SC ).
Dos 48 cursos participantes de centros universitários da região Sul, um ficou sem conceito ( SC ) e 30 receberam conceito 3, o valor modal.
Os conceitos restantes ficaram assim distribuídos : 4 ( nove cursos ), 2 ( sete cursos ) e 5 ( um curso ).
Quanto aos 24 cursos vinculados a faculdades na região Sul, 11 receberam o conceito 3, valor modal, e nove, o conceito 2.
Os 4 cursos restantes foram igualmente divididos entre os conceitos 4 e 5.
Na região Centro-Oeste, 14 dos 70 cursos oferecidos eram mantidos por universidades.
Para esses, foram atribuídos os conceitos 3 e 4 ( seis cursos cada ) e o conceito 2 ( dois cursos ).
Os centros universitários dessa região contaram com 26 cursos : 11 foram avaliados com o conceito 3 ( modal ) ; sete, com o conceito 2 ; seis, com o conceito 4 ; um, com o conceito 5 ; e um obteve o conceito 1.
Nenhum curso ficou sem conceito ( SC ).
Dos 30 cursos oferecidos em faculdades na região Centro-Oeste, 19 receberam o conceito 2 ( modal ).
Os conceitos restantes ficaram assim distribuídos : 3 ( sete cursos ) e 1 ( quatro cursos ).
Novamente, no que se refere aos CEFET/IF da região Centro-Oeste, nenhum curso foi avaliado.
De maneira semelhante à forma como se procedeu em relação à seção anterior, esta seção encerra -se com o Gráfico 4.3, que apresenta a distribuição do Conceito Enade dos cursos de Fisioterapia segundo a organização acadêmica da IES.
No Gráfico 4.3, apresenta -se a distribuição do Conceito Enade segundo a organização acadêmica da IES.
A leitura desse último gráfico é semelhante aos anteriores, seguindo -se a lógica de que poligonais deslocadas para a esquerda correspondem aos conceitos mais baixos ( 1, 2 e 3 ) e poligonais deslocadas para a direita correspondem a Conceitos Enade mais altos.
Nas universidades ( linha azul ), todos os cursos obtiveram conceito 5 ou abaixo disso, já que a poligonal cumulativa alcança 100 \% para o conceito 5.
Além disso, os cursos de universidades apresentam uma poligonal mais à direita do que as demais, denotando uma melhor distribuição de Conceitos Enade.
No outro extremo, encontram -se os cursos mantidos por faculdades ( linha vermelha ), apresentando as respectivas poligonais mais à esquerda e, consequentemente, uma distribuição de Conceitos Enade com os menores valores.
Os cursos mantidos por centro universitários ( linha verde ) apresentam poligonais mais centrais e, consequentemente, uma distribuição de Conceitos Enade com os valores intermediários.

% matched lemmas: 
\end{textsample}
